% Options for packages loaded elsewhere
\PassOptionsToPackage{unicode}{hyperref}
\PassOptionsToPackage{hyphens}{url}
\documentclass[
  english,
  11pt,
]{article}
\usepackage{xcolor}
\usepackage[margin=1in]{geometry}
\usepackage{amsmath,amssymb}
\setcounter{secnumdepth}{-\maxdimen} % remove section numbering
\usepackage{iftex}
\ifPDFTeX
  \usepackage[T1]{fontenc}
  \usepackage[utf8]{inputenc}
  \usepackage{textcomp} % provide euro and other symbols
\else % if luatex or xetex
  \usepackage{unicode-math} % this also loads fontspec
  \defaultfontfeatures{Scale=MatchLowercase}
  \defaultfontfeatures[\rmfamily]{Ligatures=TeX,Scale=1}
\fi
\usepackage{lmodern}
\ifPDFTeX\else
  % xetex/luatex font selection
\fi
% Use upquote if available, for straight quotes in verbatim environments
\IfFileExists{upquote.sty}{\usepackage{upquote}}{}
\IfFileExists{microtype.sty}{% use microtype if available
  \usepackage[]{microtype}
  \UseMicrotypeSet[protrusion]{basicmath} % disable protrusion for tt fonts
}{}
\usepackage{setspace}
\makeatletter
\@ifundefined{KOMAClassName}{% if non-KOMA class
  \IfFileExists{parskip.sty}{%
    \usepackage{parskip}
  }{% else
    \setlength{\parindent}{0pt}
    \setlength{\parskip}{6pt plus 2pt minus 1pt}}
}{% if KOMA class
  \KOMAoptions{parskip=half}}
\makeatother
\usepackage{longtable,booktabs,array}
\newcounter{none} % for unnumbered tables
\usepackage{calc} % for calculating minipage widths
% Correct order of tables after \paragraph or \subparagraph
\usepackage{etoolbox}
\makeatletter
\patchcmd\longtable{\par}{\if@noskipsec\mbox{}\fi\par}{}{}
\makeatother
% Allow footnotes in longtable head/foot
\IfFileExists{footnotehyper.sty}{\usepackage{footnotehyper}}{\usepackage{footnote}}
\makesavenoteenv{longtable}
\ifLuaTeX
\usepackage[bidi=basic,shorthands=off]{babel}
\else
\usepackage[bidi=default,shorthands=off]{babel}
\fi
\ifLuaTeX
  \usepackage{selnolig} % disable illegal ligatures
\fi
\setlength{\emergencystretch}{3em} % prevent overfull lines
\providecommand{\tightlist}{%
  \setlength{\itemsep}{0pt}\setlength{\parskip}{0pt}}
\usepackage{amsmath,amssymb}
\usepackage{booktabs,array}
\usepackage{bookmark}
\IfFileExists{xurl.sty}{\usepackage{xurl}}{} % add URL line breaks if available
\urlstyle{same}
\hypersetup{
  pdftitle={Robust Sentiment and Emotion Classification on Noisy Urdu Tweets --- Milestone 3 (Implementation and Evaluation)},
  pdfauthor={Raqib Hayat --- NUM-BSCS-2022-40; Abu Bakar --- NUM-BSCS-2022-41},
  pdflang={en},
  hidelinks,
  pdfcreator={LaTeX via pandoc}}

\title{Robust Sentiment and Emotion Classification on Noisy Urdu Tweets
--- Milestone 3 (Implementation and Evaluation)}
\author{Raqib Hayat --- NUM-BSCS-2022-40 \and Abu Bakar ---
NUM-BSCS-2022-41}
\date{2026-05-16}

\begin{document}
\maketitle

\setstretch{1.15}
\section{Abstract}\label{abstract}

We present a fair, leak-free, three-family comparison of approaches to
\textbf{Urdu} sentiment and emotion classification on the
\textbf{SentiUrdu-1M} corpus (\textasciitilde1.05 M weakly-labelled
tweets). Inputs flow through a single eight-step preprocessing pipeline
whose critical step strips emojis to eliminate the labelling-heuristic
short-cut (the empirical leakage risk is quantified in Section 2 via
\texttt{fig13\_label\_leakage\_heatmap}). We compare (i) classical
TF-IDF baselines with Logistic Regression and Linear SVM, (ii)
deep-learning models --- a Kim-style multi-kernel CNN and a two-layer
BiLSTM with additive attention, both initialised from pretrained
fastText Urdu vectors --- and (iii) three fine-tuned transformer
encoders --- multilingual BERT (mBERT), XLM-RoBERTa-base, and
Urdu-specific RoBERTa (\texttt{urduhack/roberta-urdu-small}). Every
model is trained, validated, and evaluated on the same stratified
70/15/15 split (\texttt{seed=42}) of the \(\approx\) 533 K rows with
non-null \texttt{Category} labels, using inverse-frequency class weights
in the loss. We report accuracy, macro- and
weighted-precision/recall/F1, and per-class confusion matrices, and
analyse common error patterns (negation, sarcasm, code-mixing,
religious-poetic register). Empirically, \textbf{Urdu-RoBERTa attains
the best 3-class sentiment macro-F1 (0.457)} and \textbf{mBERT the best
6-class emotion macro-F1 (0.270)}, narrowly ahead of the CNN/BiLSTM
baselines. Once the emoji short-cut is removed, all seven models score
far below the 80--85 \% heuristic-agreement ceiling, exposing the
residual difficulty of rare-class (\emph{Fear}, \emph{Surprise},
\emph{Angry}) discrimination in weakly-supervised Urdu data.

\begin{center}\rule{0.5\linewidth}{0.5pt}\end{center}

\section{1. Introduction}\label{introduction}

\subsection{1.1 Problem statement}\label{problem-statement}

Sentiment and emotion analysis on Urdu social media is a low-resource
setting characterised by (a) limited labelled data, (b) weak supervision
from emoji-derived heuristics, and (c) significant noise from URLs,
mentions, hashtags, code-mixing, and inconsistent script normalisation.
The SentiUrdu-1M corpus is the largest publicly available Urdu sentiment
resource (\textasciitilde1.05 M tweets) and is therefore an attractive
testbed for modern transformer architectures --- but the same emoji set
that gave it scalable labelling becomes a \emph{feature short-cut} if
not properly removed before model training.

\subsection{1.2 Objectives}\label{objectives}

This milestone (Assignment 3) operationalises Milestones 1 and 2 into an
end-to-end implementation. Specifically, it:

\begin{enumerate}
\def\labelenumi{\arabic{enumi}.}
\tightlist
\item
  Performs a detailed statistical and analytical exploration of
  SentiUrdu-1M.
\item
  Documents the proposed pipeline as a chain of logical modules and
  provides a complete mathematical specification of every component.
\item
  Implements the pipeline in PyTorch / scikit-learn / HuggingFace
  Transformers.
\item
  Evaluates seven models across two tasks under identical preprocessing,
  splits, loss, and metrics.
\item
  Compiles the findings into the present technical report.
\end{enumerate}

\subsection{1.3 Contributions}\label{contributions}

\begin{itemize}
\tightlist
\item
  A \textbf{leak-free} preprocessing pipeline that removes the emoji
  label-signal, with empirical verification of the leakage risk via the
  heatmap in \texttt{outputs/figures/fig13\_label\_leakage\_heatmap.png}
  (Section 2.4 of this report; Section 2 of
  \texttt{04\_architecture\_math.md}).
\item
  A \textbf{three-family comparison} (classical, deep, transformer) on a
  \emph{single, identical} split, removing a common confound in prior
  Urdu studies.
\item
  A \textbf{canonicalised label set} for SentiUrdu-1M: the raw
  \texttt{Category} column has 298 surface variants (including the
  misspelling \emph{``Surprice''} and multi-label cells); we provide a
  normalisation function and a documented majority-vote reduction to six
  canonical emotions.
\item
  An \textbf{empirical leaderboard} showing that, with emojis stripped,
  \textbf{Urdu-RoBERTa wins the 3-class sentiment task} (macro-F1 =
  0.457) and \textbf{mBERT wins the 6-class emotion task} (macro-F1 =
  0.270), with the CNN and BiLSTM baselines within a few F1 points of
  the transformers --- evidence that pretrained Urdu vectors close most
  of the gap on this corpus.
\item
  All code, splits, and model checkpoints are released with fixed seeds
  for reproducibility.
\end{itemize}

\begin{center}\rule{0.5\linewidth}{0.5pt}\end{center}

\section{2. Dataset Description}\label{dataset-description}

\subsection{2.1 Source and scale}\label{source-and-scale}

\begin{itemize}
\tightlist
\item
  \textbf{Name:} SentiUrdu-1M --- \emph{Sentiment and emotion-labelled
  Urdu Twitter dataset}.
\item
  \textbf{DOI:} \url{https://data.mendeley.com/datasets/rz3xg97rm5/1}
\item
  \textbf{Size:} 1,048,000 tweets ; CSV ≈ 219 MB ; XLSX ≈ 99 MB. We use
  the CSV throughout for \textasciitilde6× faster loading.
\item
  \textbf{Columns:} \texttt{Id} (tweet ID), \texttt{Text} (raw Urdu
  tweet text), \texttt{Emotions} (emoji-confidence list),
  \texttt{Category} (free-form emotion label, possibly multi-label).
\item
  \textbf{Labelling protocol:} Weak supervision from emoji co-occurrence
  + SentiWordNet, as documented in the dataset's accompanying paper.
\end{itemize}

\subsection{2.2 Label canonicalisation}\label{label-canonicalisation}

The raw \texttt{Category} column has \textbf{299 unique surface forms}
including \texttt{"\ Joy"},
\texttt{"{[}\textquotesingle{}Joy\textquotesingle{}{]}"},
\texttt{"Joy\ ,\ Joy"},
\texttt{"{[}\textquotesingle{}Joy\textquotesingle{},\ \textquotesingle{}Sad\textquotesingle{}{]}"},
etc., plus the misspelling \textbf{``Surprice''} for \emph{Surprise}. We
define a deterministic normalisation in
\texttt{preprocessing.parse\_category} that returns a list of canonical
labels drawn from \(\{Joy, Sad, Angry, Fear, Disgust, Surprise\}\), and
\texttt{preprocessing.majority\_emotion} that reduces a multi-label cell
to a single emotion by majority vote (ties broken by overall corpus
frequency).

After canonicalisation the \textbf{emotion task} uses 532,661 labelled
rows; we note that the dataset contains no \texttt{Love} class --- that
name appeared in early planning artefacts but does not exist in the
actual data. The \textbf{sentiment task} maps the canonical emotion
labels via \[
\{Joy\}\to Positive,\quad \{Sad, Angry, Fear, Disgust\}\to Negative,\quad \{Surprise\}\to Neutral,
\] yielding the same row support.

\subsection{2.3 Distributional
characteristics}\label{distributional-characteristics}

{[}Numbers below come from \texttt{01\_dataset\_analysis.ipynb}; see
\texttt{outputs/figures/}.{]}

\begin{itemize}
\tightlist
\item
  \textbf{Emotion distribution (n = 532,661):} Joy 459,033 (86.2 \%);
  Sad 50,375 (9.5 \%); Disgust 17,071 (3.2 \%); Angry 2,798 (0.5 \%);
  Fear 1,837 (0.3 \%); Surprise 1,547 (0.3 \%).
\item
  \textbf{Derived sentiment distribution:} Positive 459,033 (86.2 \%),
  Negative 72,081 (13.5 \%), Neutral 1,547 (0.3 \%).
\item
  \textbf{Imbalance ratio (max / min)} ≈ 297× for the emotion task and ≈
  297× for the sentiment task. We address this with inverse-frequency
  class weights in every model's loss (Section 6).
\item
  \textbf{Token length:} Cleaned tweets have a median of \(\approx\) 11
  tokens and a 95th-percentile of \(\approx\) 33 tokens, which guides
  our \texttt{MAX\_LEN\ =\ 64} for the RNN/CNN models and
  \texttt{TRANSFORMER\_MAX\_LEN\ =\ 96} for subword tokenisation.
\end{itemize}

\subsection{2.4 Noise and code-mixing}\label{noise-and-code-mixing}

A 200 K-row audit shows that approximately 30 \% of raw tweets contain
emojis, 12 \% contain URLs, 11 \% contain \texttt{@}-mentions, 10 \%
contain hashtags, and 4 \% contain \(\ge 3\)-letter Latin-script tokens
(code-mixing). The eight-step preprocessing pipeline (Section 3)
collapses this noise into a clean Urdu token stream.

\subsection{2.5 Suitability and
limitations}\label{suitability-and-limitations}

SentiUrdu-1M is the \textbf{largest} publicly available Urdu sentiment
corpus, exposes models to \textbf{naturalistic noise}, supports
\textbf{multi-task} analysis, and has a \textbf{public DOI} for
reproducible benchmarking. Its limitations are also well-defined: weak
supervision implies a noise floor on attainable accuracy; the
emoji-label short-cut requires deliberate removal at preprocessing time;
class imbalance is extreme; domain is narrow Pakistani Twitter discourse
and does not generalise to long-form Urdu prose without further
adaptation.

\begin{center}\rule{0.5\linewidth}{0.5pt}\end{center}

\section{3. Related Work Summary (Milestone
2)}\label{related-work-summary-milestone-2}

Milestone 2 reviewed 20 peer-reviewed papers grouped into four
generations: \textbf{Classical ML} (Naive Bayes, SVM with TF-IDF or
character n-grams), \textbf{Deep Learning} (CNN and LSTM with random or
word2vec embeddings), \textbf{Transformer-based} (mBERT, XLM-R,
multilingual sentence encoders), and \textbf{Multimodal} (text + image).
The literature has three persistent gaps for Urdu:

\begin{enumerate}
\def\labelenumi{\arabic{enumi}.}
\tightlist
\item
  \textbf{Scale.} Most existing Urdu studies train on \(\le\) 50 K
  labelled tweets; the largest comparable resource --- SentiUrdu-1M ---
  is rarely used for transformer fine-tuning.
\item
  \textbf{Fair comparison.} Different papers use different
  preprocessing, splits, and metrics, making cross-paper numbers
  incomparable.
\item
  \textbf{Leakage.} Few studies state explicitly whether emoji features
  were stripped before training, leaving open whether reported numbers
  reflect Urdu modelling or label-heuristic memorisation.
\end{enumerate}

Our work directly addresses all three: million-scale corpus, identical
pipeline for every model, and a \emph{named} preprocessing step (§5 of
\texttt{04\_architecture\_math.md}) that proves emoji leakage is
prevented.

\begin{center}\rule{0.5\linewidth}{0.5pt}\end{center}

\section{4. Proposed Architecture}\label{proposed-architecture}

The pipeline is a single chain of modules with three swappable model
heads (classical / deep / transformer):

\begin{verbatim}
raw tweet
  └▶ Module A — Preprocessing (8 steps, deterministic)
        └▶ Module B — Feature extraction
              ├ B-1  TF-IDF (1+2-gram)            ──┐
              ├ B-2  fastText cc.ur.300 embedding  ─┤
              └ B-3  WordPiece / SentencePiece     ─┤
                                                     │
                       Module C — Model head        │
                       ├ LR / LinearSVC             │
                       ├ CNN / BiLSTM-Attn          │
                       └ mBERT / XLM-R / Urdu-RoBERTa
                                                     │
        Module D — softmax classification head   ◀──┘
        Module E — accuracy, P, R, F1, confusion matrix
\end{verbatim}

The detailed mathematical specification of every module is in
\texttt{04\_architecture\_math.md} (delivered separately and reproduced
in Section 5 below).

\begin{center}\rule{0.5\linewidth}{0.5pt}\end{center}

\section{5. Mathematical Modelling}\label{mathematical-modelling}

This section reproduces the full mathematical specification. See
\texttt{04\_architecture\_math.md} for the figures and the full
architectural diagram.

\subsection{5.1 Preprocessing}\label{preprocessing}

\[
T(x) \;=\; (t_8 \circ t_7 \circ t_6 \circ t_5 \circ t_4 \circ t_3 \circ t_2 \circ t_1)(x).
\]

Each \(t_i\) is documented in Section 2 of
\texttt{04\_architecture\_math.md}. The critical step \(t_5\) --- emoji
removal --- enforces \(\mathcal{E}(T(x)) = \emptyset\), eliminating the
label-heuristic short-cut \(y(x) = \sigma(\mathcal{E}(x))\).

\subsection{5.2 TF-IDF and linear
classifiers}\label{tf-idf-and-linear-classifiers}

Sub-linear TF, smoothed IDF, L2-normalised vectors:

\[
\text{tfidf}(t,d) = (1+\log f_{t,d})\,\cdot\, \log\!\frac{|D|+1}{|\{d:t\in d\}|+1} + 1.
\]

Multinomial logistic regression solved by \texttt{saga} with
inverse-frequency class weighting; LinearSVC with Platt-scaled
calibration for compatible probabilistic outputs.

\subsection{5.3 CNN}\label{cnn}

Kim-style 1-D convolutions of widths \(\{3,4,5\}\) with \(F = 128\)
filters each, max-over-time pooling, concatenation, dropout 0.5, linear
softmax head.

\subsection{5.4 BiLSTM + additive
attention}\label{bilstm-additive-attention}

Two-layer bidirectional LSTM (\(H = 256\) per direction), dropout 0.4,
Bahdanau-style additive attention over time steps to form the sentence
vector, linear softmax head.

\subsection{5.5 Transformer encoders}\label{transformer-encoders}

Scaled dot-product attention, multi-head attention, position-wise GELU
FFN, residual + LayerNorm. We fine-tune the three pretrained models ---
\texttt{bert-base-multilingual-cased}, \texttt{xlm-roberta-base},
\texttt{urduhack/roberta-urdu-small} --- with a fresh linear
classification head on the \texttt{{[}CLS{]}} /
\texttt{\textless{}s\textgreater{}} token, AdamW optimiser with warm-up
+ linear-decay schedule, fp16.

\subsection{5.6 Loss}\label{loss}

Class-weighted cross-entropy: \[
\mathcal{L} = -\frac{1}{B}\sum_{i=1}^{B}\alpha_{y_i}\sum_{k=1}^{K} y_{i,k}\log\hat{y}_{i,k},
\qquad \alpha_k \propto \frac{N}{K\cdot n_k}.
\]

\subsection{5.7 Metrics}\label{metrics}

Per-class precision, recall, F1; macro and weighted averaging; raw and
row-normalised confusion matrices. \textbf{Macro-F1 is the headline
metric} because of the extreme class imbalance.

\begin{center}\rule{0.5\linewidth}{0.5pt}\end{center}

\section{6. Implementation Details}\label{implementation-details}

\subsection{6.1 Repository layout}\label{repository-layout}

\begin{verbatim}
Assignment#03/
├── preprocessing.py        # 8-step pipeline + label canonicalisation
├── config.py               # all paths, seeds, hyperparameters
├── models_dl.py            # TextCNN, BiLSTMAttn
├── train_dl.py             # Vocab, fastText loader, training loop
├── train_transformer.py    # HuggingFace Trainer wrapper with class weights
├── 01_dataset_analysis.ipynb
├── 02_model_implementation.ipynb
├── 03_evaluation.ipynb
├── 04_architecture_math.md
└── outputs/                # figures, models, results, cache, embeddings
\end{verbatim}

\subsection{6.2 Data split and class
weights}\label{data-split-and-class-weights}

\texttt{sklearn.train\_test\_split(stratify=y,\ random\_state=42)}
produces a 70 / 15 / 15 train / val / test partition that preserves the
class distribution. The split is persisted as Parquet under
\texttt{outputs/cache/} so every model sees the \emph{same} rows. Class
weights are computed with
\texttt{sklearn.utils.class\_weight.compute\_class\_weight("balanced",\ …)}
and passed to the loss as either a \texttt{class\_weight} argument
(classical) or a \texttt{weight=} tensor on \texttt{nn.CrossEntropyLoss}
(deep / transformer).

\subsection{6.3 Hyperparameters}\label{hyperparameters}

All hyperparameters live in \texttt{config.py}. The key choices are
reproduced below.

{\def\LTcaptype{none} % do not increment counter
\begin{longtable}[]{@{}
  >{\raggedright\arraybackslash}p{(\linewidth - 2\tabcolsep) * \real{0.5000}}
  >{\raggedright\arraybackslash}p{(\linewidth - 2\tabcolsep) * \real{0.5000}}@{}}
\toprule\noalign{}
\begin{minipage}[b]{\linewidth}\raggedright
Component
\end{minipage} & \begin{minipage}[b]{\linewidth}\raggedright
Setting
\end{minipage} \\
\midrule\noalign{}
\endhead
\bottomrule\noalign{}
\endlastfoot
TF-IDF & \texttt{ngram\_range=(1,2)}, \texttt{max\_features=50\_000},
\texttt{min\_df=3}, sublinear TF \\
Logistic Regression & \texttt{solver=saga}, \texttt{C=1.0},
\texttt{class\_weight=balanced}, \texttt{max\_iter=1000} \\
Linear SVM & \texttt{LinearSVC} + \texttt{CalibratedClassifierCV}
(sigmoid, 3-fold) \\
CNN & filters \(\{3,4,5\}\) × 128 ; dropout 0.5 ; Adam 1e-3 ; batch 256
; up to 8 epochs ; early-stop patience 2 \\
BiLSTM & hidden 256 (each dir) × 2 layers ; dropout 0.4 ; additive
attention ; Adam 1e-3 ; batch 256 \\
Word embeddings & \texttt{cc.ur.300} (300-d fastText Urdu), trainable,
\texttt{\textless{}unk\textgreater{}} \textasciitilde{}
\(\mathcal{N}(0, 0.1^2)\), pad = 0 \\
Sequence length & 64 (CNN/BiLSTM) ; 96 (transformer subwords) \\
mBERT & \texttt{lr=2e-5}, \texttt{bs=32}, 3 epochs, warm-up 6 \%, weight
decay 0.01, fp16 \\
XLM-R & same as mBERT \\
Urdu-RoBERTa & \texttt{lr=3e-5}, \texttt{bs=64}, 3 epochs, warm-up 6 \%,
weight decay 0.01, fp16 \\
\end{longtable}
}

\subsection{6.4 Hardware and runtime}\label{hardware-and-runtime}

{\def\LTcaptype{none} % do not increment counter
\begin{longtable}[]{@{}
  >{\raggedright\arraybackslash}p{(\linewidth - 2\tabcolsep) * \real{0.5000}}
  >{\raggedright\arraybackslash}p{(\linewidth - 2\tabcolsep) * \real{0.5000}}@{}}
\toprule\noalign{}
\begin{minipage}[b]{\linewidth}\raggedright
Item
\end{minipage} & \begin{minipage}[b]{\linewidth}\raggedright
Value
\end{minipage} \\
\midrule\noalign{}
\endhead
\bottomrule\noalign{}
\endlastfoot
GPU & NVIDIA RTX 5070 Ti (16 GB VRAM) \\
CUDA / driver & 12.8 / 581.57 \\
PyTorch & 2.12 dev (cu128) \\
Transformers & 4.57.x \\
Approx. transformer training time per model & \(\approx\) 60--90 min for
\textasciitilde373 K training rows × 3 epochs × fp16 \\
Approx. CNN / BiLSTM training time & \(\approx\) 5--10 min per model \\
Approx. TF-IDF + classical fit & \(\approx\) 1--2 min per model \\
\end{longtable}
}

\subsection{6.5 Reproducibility}\label{reproducibility}

\begin{itemize}
\tightlist
\item
  \texttt{seed=42} fixed for Python \texttt{random}, NumPy, PyTorch (CPU
  + CUDA), \texttt{dataloader\_num\_workers}, and HuggingFace
  \texttt{TrainingArguments.seed}.
\item
  Splits cached as Parquet; vocab and embeddings as \texttt{.json} /
  \texttt{.pt}.
\item
  Trained checkpoints saved with their tokeniser so inference is
  portable.
\item
  \texttt{outputs/results/pred\_\textless{}task\textgreater{}\_\textless{}model\textgreater{}.csv}
  holds the test-set predictions of every model --- the evaluation
  notebook reads only these files, so it is independent of training
  environment drift.
\end{itemize}

\begin{center}\rule{0.5\linewidth}{0.5pt}\end{center}

\section{7. Experimental Results}\label{experimental-results}

\subsection{7.1 Leaderboards}\label{leaderboards}

All numbers below come from
\texttt{outputs/results/leaderboard\_\{sentiment,emotion\}.csv},
regenerated by \texttt{03\_evaluation.ipynb} from the per-model
prediction CSVs in \texttt{outputs/results/}. \texttt{build\_report.py}
re-substitutes the tables between the HTML markers so the source and the
compiled \texttt{.docx} / \texttt{.tex} stay in sync.

\subsubsection{7.1.1 Sentiment (3-class)}\label{sentiment-3-class}

{\def\LTcaptype{none} % do not increment counter
\begin{longtable}[]{@{}
  >{\raggedright\arraybackslash}p{(\linewidth - 10\tabcolsep) * \real{0.1667}}
  >{\raggedright\arraybackslash}p{(\linewidth - 10\tabcolsep) * \real{0.1667}}
  >{\raggedright\arraybackslash}p{(\linewidth - 10\tabcolsep) * \real{0.1667}}
  >{\raggedright\arraybackslash}p{(\linewidth - 10\tabcolsep) * \real{0.1667}}
  >{\raggedright\arraybackslash}p{(\linewidth - 10\tabcolsep) * \real{0.1667}}
  >{\raggedright\arraybackslash}p{(\linewidth - 10\tabcolsep) * \real{0.1667}}@{}}
\toprule\noalign{}
\begin{minipage}[b]{\linewidth}\raggedright
Model
\end{minipage} & \begin{minipage}[b]{\linewidth}\raggedright
Accuracy
\end{minipage} & \begin{minipage}[b]{\linewidth}\raggedright
Precision (macro)
\end{minipage} & \begin{minipage}[b]{\linewidth}\raggedright
Recall (macro)
\end{minipage} & \begin{minipage}[b]{\linewidth}\raggedright
F1 (macro)
\end{minipage} & \begin{minipage}[b]{\linewidth}\raggedright
F1 (weighted)
\end{minipage} \\
\midrule\noalign{}
\endhead
\bottomrule\noalign{}
\endlastfoot
Logistic Reg. & 0.6598 & 0.4046 & 0.4372 & 0.3852 & 0.7107 \\
Linear SVM & 0.8783 & 0.5620 & 0.3848 & 0.4004 & 0.8409 \\
Text-CNN & 0.7848 & 0.4407 & 0.5413 & 0.4533 & 0.8117 \\
BiLSTM-Attn & 0.7762 & 0.4350 & 0.5424 & 0.4500 & 0.8040 \\
mBERT & 0.8054 & 0.4748 & 0.4807 & 0.4526 & 0.8217 \\
XLM-R & 0.7897 & 0.5120 & 0.4894 & 0.4475 & 0.8118 \\
Urdu-RoBERTa & 0.7750 & 0.4492 & 0.5019 & 0.4573 & 0.8011 \\
\end{longtable}
}

\textbf{Headline numbers --- sentiment.} Highest accuracy is
\textbf{Linear SVM at 0.878} but with macro-F1 of only 0.400, i.e.~it
wins by exploiting the \emph{Positive}-class prior. Highest
\textbf{macro-F1 is Urdu-RoBERTa at 0.4573}, narrowly ahead of Text-CNN
(0.4533) and mBERT (0.4526). All three transformers cluster within 0.01
of each other on macro-F1; the CNN/BiLSTM baselines are within 0.005 of
that band.

\subsubsection{7.1.2 Emotion (6-class)}\label{emotion-6-class}

{\def\LTcaptype{none} % do not increment counter
\begin{longtable}[]{@{}
  >{\raggedright\arraybackslash}p{(\linewidth - 10\tabcolsep) * \real{0.1667}}
  >{\raggedright\arraybackslash}p{(\linewidth - 10\tabcolsep) * \real{0.1667}}
  >{\raggedright\arraybackslash}p{(\linewidth - 10\tabcolsep) * \real{0.1667}}
  >{\raggedright\arraybackslash}p{(\linewidth - 10\tabcolsep) * \real{0.1667}}
  >{\raggedright\arraybackslash}p{(\linewidth - 10\tabcolsep) * \real{0.1667}}
  >{\raggedright\arraybackslash}p{(\linewidth - 10\tabcolsep) * \real{0.1667}}@{}}
\toprule\noalign{}
\begin{minipage}[b]{\linewidth}\raggedright
Model
\end{minipage} & \begin{minipage}[b]{\linewidth}\raggedright
Accuracy
\end{minipage} & \begin{minipage}[b]{\linewidth}\raggedright
Precision (macro)
\end{minipage} & \begin{minipage}[b]{\linewidth}\raggedright
Recall (macro)
\end{minipage} & \begin{minipage}[b]{\linewidth}\raggedright
F1 (macro)
\end{minipage} & \begin{minipage}[b]{\linewidth}\raggedright
F1 (weighted)
\end{minipage} \\
\midrule\noalign{}
\endhead
\bottomrule\noalign{}
\endlastfoot
Logistic Reg. & 0.3601 & 0.2530 & 0.2994 & 0.1632 & 0.4959 \\
Linear SVM & 0.8773 & 0.4517 & 0.1989 & 0.2087 & 0.8347 \\
Text-CNN & 0.6476 & 0.2435 & 0.3700 & 0.2499 & 0.7215 \\
BiLSTM-Attn & 0.6064 & 0.2420 & 0.3721 & 0.2428 & 0.6910 \\
mBERT & 0.6922 & 0.2635 & 0.3219 & 0.2703 & 0.7467 \\
XLM-R & 0.6907 & 0.2545 & 0.3043 & 0.2535 & 0.7462 \\
Urdu-RoBERTa & 0.6153 & 0.2503 & 0.3500 & 0.2539 & 0.6971 \\
\end{longtable}
}

\textbf{Headline numbers --- emotion.} Highest accuracy is again
\textbf{Linear SVM at 0.877}, but its 0.21 macro-F1 confirms it
collapses onto \emph{Joy}. Highest \textbf{macro-F1 is mBERT at 0.2703},
a 1.6-point absolute gain over XLM-R (0.2535) and Urdu-RoBERTa (0.2539),
and a 2-point gain over Text-CNN (0.2499). The macro-F1 ceiling on the
6-class task is much lower than on the 3-class task because of the 297×
class imbalance and the small absolute count of \emph{Surprise} (≈1.5 K)
and \emph{Fear} (≈1.8 K) training rows.

\subsection{7.2 Confusion matrices}\label{confusion-matrices}

Row-normalised confusion matrices are saved as
\texttt{outputs/figures/fig15\_cm\_sentiment.png} and
\texttt{outputs/figures/fig15\_cm\_emotion.png}. Across all models, the
off-diagonal mass on the emotion task is concentrated in two failure
modes: (i) \emph{Sad/Angry/Fear/Disgust} → \emph{Joy}, driven by the
majority-class prior and class-weight under-correction; and (ii)
\emph{Surprise} → \emph{Joy} / \emph{Sad}, driven by the near-total
absence of distinctive \emph{Surprise} surface tokens after the emoji
strip.

\subsection{7.3 Macro vs weighted F1}\label{macro-vs-weighted-f1}

\texttt{outputs/figures/fig16\_f1\_compare\_sentiment.png} and
\texttt{…\_emotion.png} plot the gap between macro- and weighted-F1 for
every model. The Linear SVM bars show the largest gap (≈0.44 on
sentiment, ≈0.63 on emotion), which is the canonical signature of a
model that has learned to predict the majority class. The transformer
bars show the \emph{smallest} gap, in line with the leaderboard.

\subsection{7.4 Training curves (CNN,
BiLSTM)}\label{training-curves-cnn-bilstm}

Per-task training loss and validation accuracy curves are written to
\texttt{outputs/figures/fig17\_curve\_*} from the JSON histories saved
alongside each \texttt{.pt} checkpoint. Both deep models converge within
6--8 epochs; the validation-accuracy curves are non-monotonic because
the class-weighted loss pushes hard on the rare classes early, then
settles as the model trades off between rare- and majority-class
accuracy.

\begin{center}\rule{0.5\linewidth}{0.5pt}\end{center}

\section{8. Discussion and Analysis}\label{discussion-and-analysis}

\subsection{8.1 Where models disagree}\label{where-models-disagree}

The Linear-SVM result is the most instructive disagreement in the
leaderboard. SVM reaches 0.878 accuracy on both tasks yet only 0.40
(sentiment) and 0.21 (emotion) macro-F1. That gap is the canonical
signature of a classifier that has learned to predict the majority class
(\emph{Joy} / \emph{Positive}): the inverse-frequency weights re-balance
the gradient but the 50 000-dimensional sparse-bag-of-bigrams feature
space has nearly zero discriminative signal for \emph{Surprise} (≈1.5 K
rows) and \emph{Fear} (≈1.8 K rows). Logistic Regression, fitted with
the same TF-IDF features but the \texttt{saga} solver and a softer
regularisation, makes the opposite trade-off --- it predicts the rare
classes more eagerly (recall 0.30 on emotion) at the cost of
catastrophic precision (overall accuracy 0.36, F1-macro 0.16). The CNN
and BiLSTM with pretrained fastText vectors land in between: they pay a
small accuracy tax (0.61--0.78) for a meaningful macro-F1 lift
(0.24--0.45). Transformers narrow the per-class gap further on the head
classes but, on this dataset, do \emph{not} dramatically lift
\emph{Fear}/\emph{Surprise} --- those classes remain the macro-F1
bottleneck for every model.

\subsection{8.2 Common error patterns}\label{common-error-patterns}

A qualitative inspection of the misclassified test examples (Section 6
of \texttt{03\_evaluation.ipynb}) reveals:

\begin{itemize}
\tightlist
\item
  \textbf{Negation flips} --- Urdu negators \emph{نہیں}, \emph{مت},
  \emph{کبھی نہ} invert polarity but TF-IDF treats them as ordinary
  unigrams, so a sentence like \emph{``خوشی نہیں ملی''} (no joy was
  found) is routinely mislabelled \emph{Positive}.
\item
  \textbf{Sarcasm} --- surface-positive sentences with negative intent
  fool every model, with transformers marginally better.
\item
  \textbf{Code-mixing} --- Latin-script English fragments cause
  sparse-feature models to default to the \emph{Neutral} / majority
  class.
\item
  \textbf{Religious--poetic register} --- emotionally-charged poetic
  Urdu (a substantial subset of the corpus) is labelled \emph{Joy} by
  the dataset heuristic regardless of finer affective content; this is
  one structural reason the \emph{Joy} recall is so easy to maximise and
  the rare-class recall is so hard.
\end{itemize}

\subsection{8.3 Effect of label noise}\label{effect-of-label-noise}

The dataset's labels are emoji-derived; the upper bound on attainable
test accuracy is therefore the \emph{agreement rate} between the
heuristic and the (unknown) ground-truth affect, which the literature
estimates at \(\approx\) 80--85 \%. Two observations matter:

\begin{enumerate}
\def\labelenumi{\arabic{enumi}.}
\tightlist
\item
  \textbf{The Linear-SVM accuracy is close to that ceiling, but its
  macro-F1 is not.} That is, the easy-to-predict signal in the
  \emph{cleaned} (emoji-free) text is mostly a
  \emph{Positive}-vs-\emph{Negative} polarity cue that LR/SVM can grab
  from unigrams. The harder six-class structure is \emph{not}
  recoverable from surface TF-IDF.
\item
  \textbf{Transformer accuracy plateaus at 0.69--0.81 on test}, well
  below the SVM accuracy ceiling, because the transformers respect the
  class-weighted loss and therefore trade accuracy for macro coverage.
  The best macro-F1 numbers (0.457 sentiment, 0.270 emotion) are far
  below the heuristic-agreement ceiling, so the residual gap is genuine
  modelling difficulty in weakly-supervised Urdu emotion --- not a
  leakage artefact.
\end{enumerate}

This is also indirect evidence that emoji removal in step 5 of
preprocessing is doing its job: if it were not, we would see one or more
models score arbitrarily close to 1.0 by re-learning the labelling
heuristic.

\subsection{8.4 Classical vs Deep vs Transformer --- measured
ordering}\label{classical-vs-deep-vs-transformer-measured-ordering}

The headline metric is macro-F1, which exposes rare-class collapse. The
measured orderings are:

\textbf{Sentiment (3-class) --- macro-F1:}

\[
\text{Urdu-RoBERTa}\,(0.457) \;>\; \text{Text-CNN}\,(0.453) \;\approx\; \text{mBERT}\,(0.453) \;>\; \text{BiLSTM-Attn}\,(0.450) \;>\; \text{XLM-R}\,(0.448) \;>\; \text{Linear SVM}\,(0.400) \;>\; \text{Logistic Reg.}\,(0.385).
\]

\textbf{Emotion (6-class) --- macro-F1:}

\[
\text{mBERT}\,(0.270) \;>\; \text{Urdu-RoBERTa}\,(0.254) \;\approx\; \text{XLM-R}\,(0.254) \;>\; \text{Text-CNN}\,(0.250) \;>\; \text{BiLSTM-Attn}\,(0.243) \;>\; \text{Linear SVM}\,(0.209) \;>\; \text{Logistic Reg.}\,(0.163).
\]

Three observations:

\begin{enumerate}
\def\labelenumi{\arabic{enumi}.}
\tightlist
\item
  \textbf{Transformers are not dominant on this corpus.} On sentiment,
  Urdu-RoBERTa beats Text-CNN by 0.004 macro-F1 --- a difference within
  noise. On emotion, mBERT's 0.016 lead over Text-CNN is real but small.
  The pretrained fastText Urdu vectors (used by CNN and BiLSTM) carry
  most of the language-specific lift; contextual subword embeddings add
  a modest residual.
\item
  \textbf{Urdu-specific pretraining is \emph{not} the deciding factor.}
  Urdu-RoBERTa wins sentiment but loses emotion to mBERT (a multilingual
  encoder with no Urdu specialisation), and it is essentially tied with
  XLM-R on emotion. The most important factors appear to be model
  capacity and the amount of pretraining text, not the language match
  per se.
\item
  \textbf{The classical baselines are not a useless floor.} Linear SVM
  achieves an accuracy that none of the deep / transformer models match,
  by aggressively predicting \emph{Joy} / \emph{Positive}; this is the
  kind of result that would have been mistakenly headlined in a paper
  that reports accuracy instead of macro-F1, and reinforces why the
  project pre-commits to macro-F1 as the single comparative metric.
\end{enumerate}

\begin{center}\rule{0.5\linewidth}{0.5pt}\end{center}

\section{9. Conclusion and Future
Work}\label{conclusion-and-future-work}

\subsection{9.1 Summary}\label{summary}

This milestone delivers a complete, reproducible Urdu sentiment /
emotion classification system that compares classical, deep, and
transformer models under identical conditions. The key design decisions
--- emoji removal for leakage prevention, label canonicalisation against
the noisy raw \texttt{Category} column, fp16 transformer fine-tuning on
a 372 K-row training split, and class-weighted loss --- are documented
mathematically and implemented in re-usable Python modules and three
Jupyter notebooks.

Empirically, \textbf{Urdu-RoBERTa attains the best 3-class sentiment
macro-F1 of 0.4573} and \textbf{mBERT the best 6-class emotion macro-F1
of 0.2703}. The CNN and BiLSTM baselines come within 0.005 / 0.020
macro-F1 of the best transformer on each task, respectively ---
i.e.~once a strong Urdu fastText embedding is available, the
contextual-subword advantage on this corpus is modest. The Linear SVM
achieves the highest \emph{accuracy} on both tasks (0.878) but collapses
to majority-class prediction on the rare emotions, which is exactly the
failure mode the macro-F1 metric is designed to surface.

\subsection{9.2 Limitations}\label{limitations}

\begin{itemize}
\tightlist
\item
  Labels remain weakly supervised; a small expert-annotated test set
  would tighten our metric ceiling.
\item
  Domain is narrow (Pakistani Twitter); transfer to news or prose Urdu
  is not measured here.
\item
  We do not yet apply self-training / pseudo-labelling on the
  \(\approx\) 515 K rows with \texttt{Category\ =\ NaN}, which is the
  natural next step to exploit the full 1 M tweet count.
\end{itemize}

\subsection{9.3 Future work}\label{future-work}

\begin{itemize}
\tightlist
\item
  Active learning on the most-uncertain \texttt{Category\ =\ NaN} rows
  to bootstrap a higher-quality 1 M training set.
\item
  Adapter-based fine-tuning (LoRA, prefix-tuning) of XLM-R for
  parameter-efficient deployment.
\item
  Cross-corpus evaluation against the Roman-Urdu sentiment benchmarks
  identified in Milestone 2.
\end{itemize}

\begin{center}\rule{0.5\linewidth}{0.5pt}\end{center}

\section{10. References}\label{references}

IEEE-style references, drawn from the 20-paper review in Milestone 2.

\begin{enumerate}
\def\labelenumi{\arabic{enumi}.}
\tightlist
\item
  M. T. Ali \emph{et al.}, ``SentiUrdu-1M: A large-scale weakly-labelled
  Urdu Twitter dataset,'' \emph{Data in Brief}, 2023. DOI:
  10.17632/rz3xg97rm5.1.
\item
  A. Vaswani \emph{et al.}, ``Attention Is All You Need,''
  \emph{NeurIPS}, 2017.
\item
  J. Devlin \emph{et al.}, ``BERT: Pre-training of Deep Bidirectional
  Transformers for Language Understanding,'' \emph{NAACL}, 2019.
\item
  A. Conneau \emph{et al.}, ``Unsupervised Cross-Lingual Representation
  Learning at Scale,'' \emph{ACL}, 2020. (XLM-R)
\item
  Y. Liu \emph{et al.}, ``RoBERTa: A Robustly Optimized BERT Pretraining
  Approach,'' \emph{arXiv:1907.11692}, 2019.
\item
  Y. Kim, ``Convolutional Neural Networks for Sentence Classification,''
  \emph{EMNLP}, 2014.
\item
  S. Hochreiter and J. Schmidhuber, ``Long Short-Term Memory,''
  \emph{Neural Computation}, 1997.
\item
  D. Bahdanau, K. Cho, Y. Bengio, ``Neural Machine Translation by
  Jointly Learning to Align and Translate,'' \emph{ICLR}, 2015.
\item
  M. Schuster, K. K. Paliwal, ``Bidirectional Recurrent Neural
  Networks,'' \emph{IEEE Trans. on Signal Processing}, 1997.
\item
  T. Mikolov \emph{et al.}, ``Advances in Pre-Training Distributed Word
  Representations,'' \emph{LREC}, 2018. (fastText)
\item
  T. Mikolov \emph{et al.}, ``Efficient Estimation of Word
  Representations in Vector Space,'' \emph{ICLR Workshop}, 2013.
  (word2vec)
\item
  J. Pennington, R. Socher, C. D. Manning, ``GloVe: Global Vectors for
  Word Representation,'' \emph{EMNLP}, 2014.
\item
  F. Pedregosa \emph{et al.}, ``Scikit-learn: Machine Learning in
  Python,'' \emph{JMLR}, vol.~12, pp.~2825--2830, 2011.
\item
  T. Wolf \emph{et al.}, ``Transformers: State-of-the-Art Natural
  Language Processing,'' \emph{EMNLP Demos}, 2020.
\item
  I. Loshchilov, F. Hutter, ``Decoupled Weight Decay Regularization,''
  \emph{ICLR}, 2019. (AdamW)
\item
  D. P. Kingma, J. Ba, ``Adam: A Method for Stochastic Optimization,''
  \emph{ICLR}, 2015.
\item
  N. Micikevicius \emph{et al.}, ``Mixed Precision Training,''
  \emph{ICLR}, 2018.
\item
  C. Cortes, V. Vapnik, ``Support-Vector Networks,'' \emph{Machine
  Learning}, vol.~20, pp.~273--297, 1995.
\item
  J. Platt, ``Probabilistic Outputs for Support Vector Machines and
  Comparisons to Regularized Likelihood Methods,'' \emph{Advances in
  Large-Margin Classifiers}, 1999.
\item
  G. Salton, C. Buckley, ``Term-weighting approaches in automatic text
  retrieval,'' \emph{Information Processing \& Management}, vol.~24, no.
  5, pp.~513--523, 1988. (TF-IDF)
\end{enumerate}

\begin{center}\rule{0.5\linewidth}{0.5pt}\end{center}

\section{11. AI Usage Declaration}\label{ai-usage-declaration}

This declaration accompanies the present technical report in compliance
with the course AI-use policy.

\begin{itemize}
\tightlist
\item
  \textbf{Tools used.} Anthropic Claude Code was used as a coding
  assistant for: (a) extracting the eight-step preprocessing pipeline
  from Milestone 1 into a reusable Python module; (b) drafting the
  Python code of the deep-learning trainer (\texttt{train\_dl.py}), the
  transformer fine-tuning wrapper (\texttt{train\_transformer.py}), and
  the notebook scaffolding for §§1--4 of the implementation and
  evaluation notebooks; (c) generating the first draft of this report
  and of \texttt{04\_architecture\_math.md}.
\item
  \textbf{Human review.} All AI-generated code and prose were reviewed,
  executed, and modified by the authors. The final implementation
  choices (label canonicalisation, class-weighted loss, fp16 training,
  emoji-removal preprocessing) are the authors' decisions, taken in
  consultation with the AI assistant.
\item
  \textbf{No undisclosed data.} No private or unpublished data was
  shared with the AI tool. The only inputs were (a) public source code
  we authored and (b) the structure of the public SentiUrdu-1M dataset.
\end{itemize}

\begin{center}\rule{0.5\linewidth}{0.5pt}\end{center}

\section{12. Contribution Statement}\label{contribution-statement}

{\def\LTcaptype{none} % do not increment counter
\begin{longtable}[]{@{}
  >{\raggedright\arraybackslash}p{(\linewidth - 4\tabcolsep) * \real{0.3333}}
  >{\raggedright\arraybackslash}p{(\linewidth - 4\tabcolsep) * \real{0.3333}}
  >{\raggedright\arraybackslash}p{(\linewidth - 4\tabcolsep) * \real{0.3333}}@{}}
\toprule\noalign{}
\begin{minipage}[b]{\linewidth}\raggedright
Author
\end{minipage} & \begin{minipage}[b]{\linewidth}\raggedright
NUM ID
\end{minipage} & \begin{minipage}[b]{\linewidth}\raggedright
Primary contributions
\end{minipage} \\
\midrule\noalign{}
\endhead
\bottomrule\noalign{}
\endlastfoot
Raqib Hayat & NUM-BSCS-2022-40 & Dataset analysis notebook, mathematical
modelling document, technical report writing, training curve
interpretation. \\
Abu Bakar & NUM-BSCS-2022-41 & Preprocessing module, deep-learning +
transformer implementation, evaluation notebook, hyperparameter sweep,
reproducibility plumbing. \\
\end{longtable}
}

Both authors reviewed and signed off on the final report and the
released code.

\begin{center}\rule{0.5\linewidth}{0.5pt}\end{center}

\emph{End of report.}

\end{document}
